%% ==========================================================================
%%  Abstract
%%
%%  Scoped to what the report currently carries: the model, what transfers
%%  from the goods theory, and the structural facts.  When a parked section is
%%  promoted from working/, add its result here --- and lead with it.
%% ==========================================================================
\begin{abstract}
We study the problem of allocating indivisible \emph{chores} among agents in a
fair manner, with subsidies. Envy-free allocations of indivisible items are not
guaranteed to exist, but envy-freeness can be achieved by additionally providing
some subsidy to the agents; in the chores setting this subsidy is naturally read
as compensation paid to the agents who absorb the workload. We study the subsidy
required under \emph{negative dichotomous} valuations, i.e., among agents whose
marginal disutility for any chore is either zero or one:
$v_i(S) - v_i(S \cup \set{g}) \in \set{0,1}$ for every agent $i$, every
$S \subseteq \items$ and every $g \in \items \setminus S$. This is the exact
mirror of the dichotomous goods model of Barman et al.~\cite{BKNS22}, who prove
that there a per-agent subsidy of $0$ or $1$ always suffices, and we conjecture
that the same guarantee holds for chores.

The mirror is not a relabelling. Subsidies are one-signed --- money is always a
good --- so negating an instance negates the valuations but not the money, and
the goods theory does not transfer by symmetry. This report sets out what does
transfer and what the chore problem looks like once it has. The Halpern--Shah
characterisation of envy-freeable allocations~\cite{HS19} survives verbatim,
since it assumes nothing about the sign or the structure of the valuations, and
it delivers an \emph{integral} minimum subsidy vector here; the matching lower
bound of $n-1$ survives as well. Existing work already gives a bounded subsidy:
negative dichotomous valuations are doubly monotone under the two-sided
normalisation, so the reduction of Kawase et al.~\cite{KMSTY24} yields $n-1$ per
agent and $n(n-1)/2$ in total. The open question is whether the far stronger
$\set{0,1}$ guarantee of~\cite{BKNS22} survives the passage from goods to
chores, a factor of $n$ below that baseline.

We then isolate three structural facts about the chore problem: an arc-update
lemma reducing the entire goods/chore asymmetry to two lines; a two-tier
characterisation of the allocations admitting a $\set{0,1}$ subsidy; and a
size-shift transform exhibiting the chore problem as \emph{exactly} the goods
problem of~\cite{BKNS22} together with a cardinality-balancing requirement,
which is the precise sense in which the conjecture does not follow from that
work as a black box.
\end{abstract}
